% Source: https://artofproblemsolving.com/community/c3h1405948_geometry_problem
% Stable identity: aops-topic-1405948-q01
Let the sheet be $ABCD$ with $AB=8$, $BC=6$, and fold $A$ onto $C$. Every point of the crease is equidistant from $A$ and $C$, so the crease is the part of the perpendicular bisector of $AC$ inside the rectangle; it passes through the centre $M$, and it meets $AB$ at $X$ and $CD$ at $Y$.

\begin{asy}
usepackage("amsmath");
size(6cm);
defaultpen(fontsize(10pt));
path ra(pair P,pair u,pair v,real s=0.2){return P+s*unit(u)--P+s*unit(u)+s*unit(v)--P+s*unit(v);}
pair A=(0,0),B=(8,0),C=(8,6),D=(0,6);
pair X=(6.25,0),Y=(1.75,6),M=(4,3);
draw(A--B--C--D--cycle,linewidth(0.9));
draw(A--C,gray+dashed);
draw(X--Y,red+linewidth(0.9));
draw(ra(M,C-M,Y-M,0.35));
dot(A);dot(B);dot(C);dot(D);dot(X,red);dot(Y,red);dot(M);
label("$A$",A,SW); label("$B$",B,SE); label("$C$",C,NE); label("$D$",D,NW);
label("$X$",X,S,red); label("$Y$",Y,N,red); label("$M$",(4.8,2.85));
label("$8$",(4,0),S); label("$6$",(0,3),W);
label("$x$",(3.2,0),N); label("$8-x$",(7.1,0),N);
label("\mbox{}",(min(currentpicture,user=true)+max(currentpicture,user=true))/2,black); // keeps the text after the figure black when LaTeX typesets the labels inline
\end{asy}

Let $AX=x$. Then $XC=XA=x$ and $XB=8-x$, so in right triangle $XBC$
\[
(8-x)^2+6^2=x^2\quad\Longrightarrow\quad 100=16x\quad\Longrightarrow\quad x=\dfrac{25}{4}.
\]
The diagonal is $AC=10$, so $AM=5$, and in right triangle $AMX$
\[
MX=\sqrt{AX^2-AM^2}=\sqrt{\dfrac{625}{16}-25}=\sqrt{\dfrac{225}{16}}=\dfrac{15}{4}.
\]
By symmetry about $M$, $MY=MX$, so the crease has length $XY=2\cdot\dfrac{15}{4}=\boxed{\dfrac{15}{2}}$.
